\section{Conclusion}
This study rigorously characterized the noise profile and transient response dynamics of the AHT10 sensor through a strictly deterministic, localized data acquisition protocol. 
By explicitly disabling wireless hardware interrupts and applying linear detrending, the true analog noise floor was successfully isolated. 
A direct mathematical comparison against the 20-bit digital quantization boundaries strongly indicates that signal uncertainty is dominated by analog front-end noise, exceeding the theoretical LSB limit by factors of $56.8\times$ for temperature and $388\times$ for relative humidity. 

To condition the signal effectively while preserving temporal fidelity, a calculus-based Cost Function was deployed to numerically optimize independent Exponential Moving Average filters. 
The resulting parameters, $\alpha_T = 0.1494$ and $\alpha_{RH} = 0.2260$, demonstrate that disparate physical mechanisms—thermal inertia and moisture diffusion—mandate independent signal processing configurations. 
Furthermore, the convergence of the optimized time constant ($\tau_{RH} \approx 8.00$\,s) with the sensor's published physical hardware specifications successfully validates the employed mathematical framework, ensuring optimal preservation of both transient accuracy and steady-state precision in metrological evaluations.